% Fragmento público generado desde propietarios íntegros.
% Residencia: 52.5.2.
% Fuente: ../incoming/radion_monografia_20260827/manuscrito/sections/06_interior_helice.tex:109-201
\subsubsection{Elevación helicoidal del retorno nonádico}

\paragraph{Descomposición nonádica y cociente de memoria}

La operación elemental
\begin{equation}
n=9q_9(n)+\rho_9(n),
\qquad
H_9(n)=(\rho_9(n),q_9(n))
\label{eq:H9}
\end{equation}
satisface
\begin{equation}
H_9(n+9)=H_9(n)+(0,1).
\label{eq:H9-helix}
\end{equation}
La primera coordenada cierra; la segunda asciende. Esta identidad es la
versión aritmética mínima de un tornillo.

Para una fibra espectral,
\[
n\longmapsto(\tau_\gamma^n,n)
\]
es una hélice sobre \(S^1\). Añadir la contracción \((5/9)^n\) produce una
espiral interior. El radión participa en ambas lecturas: una unidad de fase
en la frontera y una unidad orientada de acción en el lift.

\paragraph{Rango de memoria como lector dimensional}

Supongamos que una proyección visible conserva el valor de fase, pero pierde
\(d\) cociclos de memoria linealmente independientes. Cualquier lift fiel
necesita al menos \(d\) coordenadas adicionales para reconstruirlos. Esto
motiva una lectura precisa de Mecánica Dimensional:
\begin{equation}
\begin{aligned}
\text{dimensión emergente}
&=\text{rango mínimo de memoria independiente}\\
&\phantom{=}\text{no publicable en la hoja visible}.
\end{aligned}
\label{eq:dimension-memory}
\end{equation}
No se afirma que toda dimensión física sea un contador. Se afirma que, en el
dominio TPK, cada memoria independiente que debe sobrevivir fuerza una
coordenada del lift.

\begin{narrativebox}
La dimensión aparece cuando la proyección ya no puede conservar toda la
historia. El círculo es suficiente para decir dónde está la fase; no es
suficiente para decir cuántas veces volvió, qué hoja transportó, qué intervalo
se contrajo o qué acarreo sobrevivió.
\end{narrativebox}

\paragraph{Ensamblaje de cilindros, espirales y trayectorias helicoidales}

Los cilindros coinductivos no son tubos físicos añadidos a una espiral. Son
intervalos encajados que conservan prefijos y estrechan el valor. La espiral
describe contracción transversal; la hélice describe memoria longitudinal;
el cilindro describe el conjunto de valores compatibles a una profundidad
finita. En el límite, la intersección de cilindros selecciona una sección
única.

La imagen espacial reúne, sin confundir:
\[
\begin{array}{ccl}
\text{ángulo} &\leftrightarrow& \text{fase visible},\\
\text{altura} &\leftrightarrow& \text{retorno/memoria},\\
\text{radio de cilindro} &\leftrightarrow& \text{incertidumbre de refinamiento},\\
\text{hoja} &\leftrightarrow& \text{orientación bidireccional}.
\end{array}
\]
La tasa \(3^{-54}\) contrae el cilindro por retorno; no se suma al ángulo y
no reemplaza el acarreo real \(\epsR\).

\subsubsection{Cierres de fase de órdenes \(2\pi\) y \(4\pi\)}

En la proyección circular, \(2\pi\) restaura la fase. En un lift orientado
con dos hojas, una vuelta puede cerrar en la hoja conjugada y necesitar una
segunda vuelta para restaurar la orientación:
\[
U_{2\pi}=-I,\qquad U_{4\pi}=I.
\]
La estructura coincide con la periodicidad espinorial cuando se aplica el
funtor a una representación \(SU(2)\); no se identifica el lift TPK con
\(SU(2)\) sin ese mapa.

La lectura areal es compatible:
\[
\mathcal A(2\pi)=\hfull,\qquad
\mathcal A(4\pi)=2\hfull.
\]
En una banda de Möbius, un circuito visible invierte hoja y dos circuitos
restauran la orientación. El \(2\pi/4\pi\) no es una duplicación accidental;
es la diferencia entre retorno de fase y retorno del estado enriquecido.
